\documentclass[UTF8,11pt]{ctexart}

\usepackage[margin=1in]{geometry}
\usepackage{amsmath,amssymb}
\usepackage{booktabs}
\usepackage{array}
\usepackage{longtable}
\usepackage{ragged2e}
\usepackage{tabularx}
\usepackage{graphicx}
\usepackage{url}
\usepackage{cite}
\usepackage{xcolor}
\usepackage{tikz}
\usepackage{hyperref}

\setlength{\emergencystretch}{2em}

\newcolumntype{L}[1]{>{\RaggedRight\arraybackslash}p{#1}}
\newcolumntype{Y}{>{\RaggedRight\arraybackslash}X}
\newcommand{\TableTight}{\footnotesize\setlength{\tabcolsep}{2pt}\renewcommand{\arraystretch}{1.10}}

\newenvironment{algorithm}[1][]{\begin{figure}[#1]}{\end{figure}}
\newenvironment{algorithmic}[1][]{\begin{enumerate}}{\end{enumerate}}
\newcommand{\Require}{\item[\textbf{输入：}]}
\newcommand{\Ensure}{\item[\textbf{输出：}]}
\newcommand{\State}{\item}
\newcommand{\For}[1]{\item[\textbf{对于}] #1 \textbf{执行}}
\newcommand{\ForAll}[1]{\item[\textbf{对于所有}] #1 \textbf{执行}}
\newcommand{\While}[1]{\item[\textbf{当}] #1 \textbf{时}}
\newcommand{\If}[1]{\item[\textbf{如果}] #1 \textbf{则}}
\newcommand{\Else}{\item[\textbf{否则}]}
\newcommand{\Return}{\textbf{返回} }
\let\REQUIRE\Require
\let\ENSURE\Ensure
\let\STATE\State
\let\FOR\For
\let\FORALL\ForAll
\let\WHILE\While
\let\IF\If
\let\ELSE\Else
\let\RETURN\Return
\newcommand{\EndFor}{}
\newcommand{\EndIf}{}
\newcommand{\EndWhile}{}
\let\ENDFOR\EndFor
\let\ENDIF\EndIf
\let\ENDWHILE\EndWhile
\newcommand{\InputTranslatedOrPending}[2]{%
  \IfFileExists{#1.tex}{\input{#1}}{%
    \section*{#2（中文翻译待完成）}%
    \addcontentsline{toc}{section}{#2（中文翻译待完成）}%
    本章完整中文翻译尚未生成。当前翻译任务被 Claude 侧速率限制和上游错误阻断；英文原文见 \texttt{../#1.tex}。%
  }%
}

\hypersetup{
  colorlinks=true,
  linkcolor=blue!60!black,
  citecolor=green!45!black,
  urlcolor=blue!60!black,
  pdftitle={自进化 AI 智能体综述：反馈驱动的生成、评估、记忆与自我修改},
  pdfauthor={aha team}
}

\title{自进化 AI 智能体综述：反馈驱动的生成、评估、记忆与自我修改}
\author{aha team}
\date{2026 年 5 月}

\begin{document}
\maketitle

\begin{abstract}
自进化 AI 智能体通过反馈驱动的改进环路，修改自身运行中的持久组件，包括提示词、记忆、工具、代码、工作流、数据课程，有时也包括模型行为。本文面向语言模型自我改进、进化计算、AutoML、智能体架构搜索、代码修复、基准测试框架和生产级智能体框架，建立统一的综述框架。我们将文献组织为五类进化环路，比较代表性论文和开源系统，总结基准证据，分析用户与生产痛点，并识别能够改进自身行动机制的智能体所需要的安全与治理要求。本文的中心论点是：自进化必须被评估为一个受控系统过程。每一个改进主张都必须说明发生了什么变化、哪类反馈支持了变化、哪个评估器批准了变化、它带来了多少成本、引入了什么风险，以及它如何被审计或回滚。
\end{abstract}

\noindent\textbf{关键词：} 自进化 AI；LLM 智能体；递归自我改进；智能体评估；记忆；代码进化；自动化智能体设计；AI 安全

\tableofcontents
\newpage

\InputTranslatedOrPending{ch1-intro}{第 1 章 引言}
\InputTranslatedOrPending{ch2-taxonomy}{第 2 章 分类体系}
\InputTranslatedOrPending{ch3-methods}{第 3 章 方法}
\InputTranslatedOrPending{ch4-evolutionary}{第 4 章 演化系统}
\InputTranslatedOrPending{ch5-evaluation}{第 5 章 评估}
\InputTranslatedOrPending{ch6-frameworks}{第 6 章 框架}
\InputTranslatedOrPending{ch7-painpoints}{第 7 章 痛点}
\InputTranslatedOrPending{ch8-future}{第 8 章 未来方向}

\appendix
\InputTranslatedOrPending{appendix-en}{附录}

\bibliographystyle{plain}
\bibliography{../references,../references-aliases}

\end{document}
